% P0 -- Reader's Guide to the Relational Carrier Theory Corpus
% Companion entry-point document for Papers 1-6 and the bridge note.

\documentclass[11pt,a4paper]{article}

\usepackage{amsmath,amssymb,amsthm}
\PassOptionsToPackage{hyphens}{url}
\usepackage[colorlinks=true,linkcolor=blue,citecolor=blue,urlcolor=blue,breaklinks=true]{hyperref}
\hypersetup{
  pdftitle={Reader's Guide to the Relational Carrier Theory Corpus (P0)},
  pdfauthor={Sandro Bucciarelli},
  pdfsubject={Relational carrier theory},
}
\makeatletter
\g@addto@macro\UrlBreaks{\do\_\do\/\do\.\do\-}
\makeatother
\usepackage[margin=2.5cm]{geometry}
\usepackage{booktabs}
\usepackage{longtable}
\usepackage{array}
\usepackage{xcolor}

% Clickable GitHub links for src/*.py citations (auto-generated)
\newcommand{\ghrepo}[1]{https://github.com/TerraSignum/#1}
% \nolinkurl gives a monospace, line-breakable rendering of the
% short script-stem display text without competing with \href.
\newcommand{\pyscript}[3]{%
  \href{\ghrepo{#1}/blob/main/#2}{\nolinkurl{#3}}}
\newcommand{\repoOwn}{relational-carrier-manifest}
\newcommand{\pyown}[2]{\pyscript{\repoOwn}{#1}{#2}}

\emergencystretch=3em
\hbadness=4000
\Urlmuskip=0mu plus 1mu\relax

\title{Reader's Guide to the Relational Carrier Theory Corpus\\
       \large (P0 --- canonical synthesis of Papers 1--6 and the
       bridge note)}
\author{Sandro Bucciarelli\thanks{Independent researcher.
        Correspondence: \href{mailto:sandro.bucciarelli89@gmail.com}{sandro.bucciarelli89@gmail.com}.}}
\date{\today}

\newcommand{\Rsys}{\mathcal{R}}
\newcommand{\axi}{\alpha_{\xi}}
\newcommand{\esync}{\varepsilon_{\mathrm{sync}}^{2}}
\newcommand{\bpi}{\beta_{\pi}}
\newcommand{\DOmega}{D_{\Omega}}

\begin{document}
\maketitle

\begin{abstract}
\noindent This document is the canonical entry point to the
relational-carrier-theory corpus. It is a synthesis, not a paper:
every claim cites the load-bearing source paper, no new numerical
result or proof is introduced here. The corpus consists of six
sector papers (P1--P6) plus a bridge note; this guide states the
theory's one-sentence thesis, identifies the three load-bearing
claims, maps each paper to its role, registers the closure grammar
(EXACT / PRECISE / FACTOR2; load-bearing / downstream / supporting /
external / open), tabulates the top-tier closures, and lists the
falsification handles that would settle the framework's main
predictions within the next $\sim\!2$ years. The intended reader is
either a physicist arriving at the corpus for the first time or a
reader wanting the architecture before the technical layers.
\end{abstract}

\section*{Preface: Why relations?}

This corpus began from a simple question: what if physical
reality is not built from objects placed inside space, but
from relations that become stable enough to define space?

A difference by itself is not yet information. It becomes
information only when it can be compared, reached, and
conserved within a common structure. In this sense, relation
is logically prior to information, and geometry to relation.

The technical work collected here explores this possibility
in a reproducible way. The starting point is a finite
relational carrier: a network of weighted relations whose
structure is read through a bounded-operator basis and
reduced to a small algebraic coefficient system. The claim
tested throughout the corpus is not that the philosophical
intuition is evidence, but that the same relational carrier
can reproduce, constrain, or organise several physical
sectors: emergent geometry, field-theoretic scales,
loop-class corrections, black-hole entropy, dark-sector
phenomenology, and falsifiable cosmological handles.

The philosophical role of this preface is therefore limited.
It is not part of the proof structure. It does not introduce
new assumptions, equations, or numerical claims. The
technical claims begin in the Reader's Guide that follows
and are owned by the individual papers. Each result stands
or falls by its definitions, audits, reproducibility
scripts, and falsification criteria.

The intuition behind the construction can be stated simply:

\begin{quote}
\emph{A world may be constituted not from things, but from
reachable differences.}
\end{quote}

If a difference cannot be reached, compared, or conserved, it
is not yet physical information. If enough differences become
globally reachable while remaining locally structured, a
relational carrier can emerge. From such a carrier, geometry
is not assumed as a background stage; it is something to be
reconstructed. The Reader's Guide that follows maps the
technical content of the corpus to specific load-bearing
claims and falsification handles.

\section*{What this document is and is not}

\textbf{What this is.} A reader's guide. A map of where each
question is settled in the corpus. A registry of the closure
grammar. A pre-organised table of the load-bearing claims and
falsification handles.

\textbf{What this is not.} Not a paper. Not a derivation. Not a
declaration of completeness. Not peer-reviewed. The relational
carrier theory is a work-in-progress closure construction by an
independent researcher; the empirical and conceptual claims of
Papers 1--6 are presented as candidate closures whose external
validation is the principal next step. \emph{This guide
introduces no results of its own}: it is a navigational digest,
so every numerical value, every closure form, and every
falsification handle is reproduced from --- and attributed to ---
the technical paper that derives and owns it. The original
content lives in Papers 1--6; if a sentence in this guide and a
sentence in the cited paper disagree, the cited paper is
authoritative.

\section{One-sentence thesis and short genealogy}

\paragraph{One-sentence thesis.} Physical reality is described
as a finite relational carrier whose correlation structure
generates the four observable channels of the bounded-operator
readout --- two primary ($\Xi$-cohesion, $\Omega$-synchronisation)
and two coupled ($\Xi$-$\psi$ common, sync-floor) --- from which
the System-$\Rsys$ rational coefficient tuple
\begin{equation*}
(\gamma,\axi,\esync,\bpi,\DOmega)
\;=\;(1/10,\,9/10,\,1/20,\,15/16,\,67/80)
\end{equation*}
is read off, and from which the emergent geometry of general
relativity, the Standard-Model fermion-mass-and-mixing
spectrum, dark-sector phenomenology, and the
Bekenstein--Hawking entropy $S/A\!=\!1/4$ follow as algebraic
consequences without additional free parameters.

\paragraph{Guiding principle.} The whole construction can be
read as a single extremal principle: \emph{physical reality is
the extremal organisation of realizable distinguishability} ---
the configuration that extremalises a structural action under
finite reconstruction resources. This is not a separate
postulate but a name for the variational backbone the corpus
already carries: the seven-sector relational action of
P6~\cite{Paper6} admits a variational principle whose
stationarity conditions $\delta\mathcal S/\delta(\cdot)\!=\!0$
slave every field to a deterministic lattice fixed point, and
the emergent Einstein closure of P4~\cite{Paper4} is
per-direction Galerkin residual minimisation. The guiding
principle states what these constructions share; its content
is exactly those published variational structures, not an
additional axiom.

\paragraph{Foundational note: relation precedes information.}
The corpus's architecture instantiates a logical priority that
is worth stating explicitly. A difference $a\!\neq\!b$ is not
information in isolation: it requires a comparison relation in
which $a$ and $b$ are mutually accessible. The relational
carrier supplies that comparison structure as its primitive
datum; the System-$\Rsys$ rational tuple, the bounded-operator
channels, and all downstream observables are read off from it
\emph{after} the relational structure is in place, not before.
In this construction the name ``relational carrier theory'' is
literal: \emph{the relations come first, the differences become
information by virtue of the relations}, and emergent geometry
and physical observables are algebraic consequences of the
relational correlation structure. The empirical Lemma~B
finding that the spectral gap $\lambda_{2}$ is globally
distributed --- not localised on any defect sub-population
(P4~\cite{Paper4})~--- is the operational counterpart of this
priority: a connected relational carrier is the precondition
for any observable to be \emph{information about} anything
else.

\paragraph{One-paragraph genealogy.} The System-$\Rsys$
primitive set is itself fully determined by two integer inputs ---
the relational-spacetime dimension $d\!=\!4$ and the fermion
generation count $N_{\rm gen}\!=\!3$ --- via
$\gamma\!=\!1/(2(d{+}1)),\;
\axi\!=\!(2d{+}1)/(2(d{+}1)),\;
\esync\!=\!1/(4(d{+}1)),\;
\bpi\!=\!15/16,\;
\DOmega\!=\!1\!-\!(N_{\rm gen}d{+}1)/(d^{2}(2 N_{\rm gen}{-}1))$
(at the anchor $d{=}4, N_{\rm gen}{=}3$ the subtracted correction
term is $(N_{\rm gen}d{+}1)/(d^{2}(2 N_{\rm gen}{-}1)) =
(N_{\rm gen}d{+}1)/(d^{2}(d{+}1)) = 13/80$, the two denominators
agreeing via the algebraic identity $2 N_{\rm gen}{-}1 = d{+}1$,
so $\DOmega = 1 - 13/80 = 67/80$; the Cabibbo-class form
$(2 N_{\rm gen}{-}1)$ is canonical, derived in Paper~4). Both $d$ and $N_{\rm gen}$ are observationally fixed integer
counts of discrete-geometric features of the underlying
relational lattice, not free continuous parameters. The count
of free continuous parameters in the structural input is zero;
this is what the corpus means by ``no free parameters''. We
emphasise the qualifier: $d\!=\!4$ and $N_{\rm gen}\!=\!3$
\emph{are} empirical inputs in the sense that observation of
our world fixes their values, and the framework does not
derive these integers from a deeper structure. This corpus is
a \emph{minimal-input program}: two integer inputs and one
additional structural anchor (the carrier-action functional
form $S_{\rm UV}$), all other parameters derived. It is not a
``zero information input'' programme in the sense of a theory
that derives spacetime dimensionality and generation count
from first principles --- that remains a separate analytical
target. Readers comparing this corpus to ``fully derived''
fundamental theories should bear this distinction in mind:
RCT is parameter-free in the continuous sense, not
input-free in the discrete-geometric sense.
The five System-$\Rsys$
rationals are the algebraic backbone of all downstream
closures.

\paragraph{Derived primitive identities.}
Several closed-form identities among the primitives recur
across the corpus and are used repeatedly in the downstream
closures. For convenience of cross-paper reading they are
registered here:
\begin{itemize}
\item $\axi + \gamma = 1$ (chirality complementarity; P2
\S~5);
\item $\esync = \gamma/2$ (fluctuation--dissipation symmetry;
P2 \S~5);
\item $\DOmega = \bpi - \gamma = 15/16 - 1/10 = 67/80$
(diffusion--projection balance; the closed-form definition of
$\DOmega$ among the primitives, closed by theorem in
P6~\cite{Paper6} and the C2 constraint of P2 \S~5);
\item $2 N_{\rm gen} - 1 = d + 1 = 5$ (Cabibbo-class identity;
makes the two $\DOmega$ genealogy denominators agree and sets
the $V_{cb}$ dimensional structure);
\item $\gamma^{2} = 1/100 = \bigl(2(d{+}1)\bigr)^{-2}$ (the
universal second-order ``leakage'' scale: it sets the
$\gamma^{2}$ prefactor of $n_{s}\!=\!1-\gamma^{2}(d{+}N_{\rm gen})/2$,
the $\gamma^{2}c_{X}$ sub-leading corrections of P4~\cite{Paper4},
and the $\alpha_\xi^{2}$ vs.\ $\alpha_\xi^{2}(1{-}\gamma^{2})$
Continuum-Backreaction-Identity ambiguity);
\item $\axi^{2} = 81/100$ (the variance-reactivity / backreaction
continuum value: it recurs as the common leg-value of the
Continuum Backreaction Identity and as the within-$P_5$
$T_{00}^{\Xi}$ Symanzik asymptote in P4~\cite{Paper4});
\item $s_{\rm face} = 1/d = 1/4$ (Bekenstein--Hawking entropy
face, equivalently the inverse-dimension prefactor; P3 L7);
\item $d + N_{\rm gen} = 7$ (mass-mode dimensional primitive;
recurs in $L4\!=\!-(d{+}N_{\rm gen})^{2}/(d(d{+}1))\!=\!-49/20$,
$n_{s}\!=\!1-\gamma^{2}(d{+}N_{\rm gen})/2\!=\!193/200$, and
$Y_{p}\!=\!(d{+}N_{\rm gen})^{2}\gamma^{2}/2\!=\!49/200$);
\item $2(d{+}1) = 10 = \gamma^{-1}$ (back-channel inverse;
the leading integer in the Coleman--Gamow action
$S_{\rm Gamow}\!=\!2(d{+}1)N_{\rm gen}-(2d{+}N_{\rm gen})\!=\!19$);
\item $2d + N_{\rm gen} = 11$ (CKM dimensional primitive;
$V_{cb}\!=\!\axi/(2(2d{+}N_{\rm gen}))\!=\!9/220$);
\item $(d{+}1)/d = 5/4$ (gravitational-backreaction
screening amplitude; cosmological-constant Schwinger--Dyson
layer);
\item $N_{\rm gen}\cdot d = 12$ (occupancy-sector
dimensional product; pairs-per-cell on the canonical
ladder).
\end{itemize}
Each identity is a closed-form consequence of the integer
inputs $(d, N_{\rm gen})$; they are not independent free
parameters.

% SANDBOX-MOVED 2026-05-16: cross-closure-identities paragraph removed
% from this manifest pending further physical-foundation analysis.
% Content preserved at
%   ../_sandbox_cross_identities/p0_cross_closure_identities_paragraph.tex
% and ../_sandbox_cross_identities/cross_identity_search_v1.py
% Reasoning: the 14 collected (4,3)-specific cross-closure rational
% identities provide a striking algebraic-web cross-check but their
% physical-foundation status is not yet rigorously established beyond
% the (d=4, N_gen=3) numeric anchoring. Postpone integration until
% the underlying common-source mechanism is articulated, to avoid
% over-claiming structural depth from anchoring-specific arithmetic.
\iffalse
\paragraph{Cross-closure algebraic identities at $(d, N_{\rm gen})\!=\!(4, 3)$.}
\label{par:cross_closure_identities}
Beyond the primitive--primitive relations above, the
$(d, N_{\rm gen})\!=\!(4, 3)$ anchoring point realises a dense
web of \emph{cross-closure} algebraic identities: independent
closures across the cosmological, electroweak, graph-statistical
and curvature sectors satisfy small-integer rational relations
that are valid at the canonical input but not in general
$(d, N_{\rm gen})$. Four of these are particularly load-bearing:
\begin{itemize}
\item $\gamma/(\Lambda_{\rm lat}^{\infty}\!-\!1)\;=\;\lambda_{\infty}
\;=\;3/8$ \quad
(universal-leakage primitive
$\gamma$~\cite{Paper2}, lattice cosmological-constant
asymptote $\Lambda_{\rm lat}^{\infty}\!=\!19/15$~\cite{Paper4},
and $(\mathrm{SG})$-axiom spectral gap
$\lambda_{\infty}\!=\!(d{-}1)/(2d)$~\cite{Paper4}).
The carrier-defect leakage scale, the departure of the lattice
cosmological constant from unity, and the spectral gap of the
weighted Laplacian are interlocked at the algebraic level: any
two pin the third.
\item $H_{0}/(100\,{\rm km/s/Mpc})\;-\;\lambda_{\infty}\;=\;
\gamma\,(d{-}1)\;=\;N_{\rm gen}\gamma\;=\;3/10$ \quad
(dimensionless Hubble parameter~\cite{Paper4B} and
$(\mathrm{SG})$-axiom spectral gap).
\item $H_{0}/(100\,{\rm km/s/Mpc})\;+\;\varepsilon_{\rm sync}^{2}
\;=\;\frac{2d+N_{\rm gen}(d{+}N_{\rm gen})}{2d(d{+}1)}
\;=\;\frac{P_{29}}{2d(d{+}1)}\;=\;29/40$ \quad
(Hubble parameter, fluctuation--dissipation sync-defect, and
the EW/Yukawa spine integer
$P_{29}\!=\!2d\!+\!N_{\rm gen}(d{+}N_{\rm gen})\!=\!29$
that anchors $m_{\mu}/m_{e}\!=\!6000/29$~\cite{Paper3}).
\item $\lambda_{\infty}\;=\;2\,N_{\rm gen}/2^{d}\;=\;3/8$,
equivalently $\lambda_{\infty}\!-\!N_{\rm gen}/2^{d}
\!=\!N_{\rm gen}/2^{d}\!=\!3/16$ \quad
($(\mathrm{SG})$-axiom related to the
Clifford-algebra-dimension $\mathrm{Cl}(d)\!=\!2^{d}$ via the
generation count; the same Clifford-dimension structure
appears in the
$T_{00}$ vacuum-branch $S_{4}$ closure
$1\!-\!N_{\rm gen}/2^{d}\!=\!13/16$
(par.~\ref{par:rho_carrier_axis_bundle} below cross-references
in P4~\cite{Paper4}) and in the Bakry--\'Emery curvature-
dimension cross-projection
$K^{p_{5}}_{\rm BE}\!=\!\lambda_{\infty}N_{\rm gen}(d{+}N_{\rm gen})/2^{d}
\!=\!63/128$~\cite{Paper4}).
\item $V_{us}\;+\;H_{0}/(100\,{\rm km/s/Mpc})
\;=\;\alpha_{\xi}\;=\;9/10$ \quad
(Cabibbo angle~\cite{Paper3}, dimensionless Hubble
parameter~\cite{Paper4B}, and System-$\Rsys$ primitive
$\alpha_{\xi}$: three closures from electroweak, cosmological,
and primitive-axis sectors sum to the carrier-axis weight
because $V_{us}\!=\!\alpha_{\xi}s_{\rm face}$ and
$H_{0}/(100\,{\rm km/s/Mpc})\!=\!\alpha_{\xi}s_{\rm face}N_{\rm gen}$
share the prefactor $\alpha_{\xi}s_{\rm face}$ and
$(N_{\rm gen}\!+\!1)/d\!=\!1$ at the canonical anchoring
$N_{\rm gen}\!=\!d\!-\!1$).
\item $\lambda_{\infty}\;+\;s_{\rm face}\;=\;(d{+}1)/(2d)
\;=\;5/8$ \quad
\emph{universal in $d$} (independent of $N_{\rm gen}$):
the $(\mathrm{SG})$-axiom spectral gap and the
Bekenstein--Hawking entropy face sum to the sync-cone-axis
fraction of the chirality-doubled spatial dimension. This is
a property of the integer $d$ alone and the
$\lambda_{\infty}\!=\!(d{-}1)/(2d)$ definition; it shows
that the $(\mathrm{SG})$ gap and the BH entropy face are
algebraically complementary on the chirality-doubled-spatial
denominator $2d$.
\item $H_{0}/(100\,{\rm km/s/Mpc})\;=\;\lambda_{\infty}\;+\;s_{\rm face}\;+\;\varepsilon_{\rm sync}^{2}\;=\;27/40$ \quad
(dimensionless Hubble parameter as the exact additive sum of
three independent fundamental quantities: the
$(\mathrm{SG})$-axiom spectral gap
$\lambda_{\infty}\!=\!3/8$, the Bekenstein--Hawking entropy
face $s_{\rm face}\!=\!1/4$, and the fluctuation--dissipation
sync-defect $\varepsilon_{\rm sync}^{2}\!=\!1/20$ from three
different corpus papers (P4, P3, P2 respectively). The
arithmetic is exact: $15/40\!+\!10/40\!+\!2/40\!=\!27/40$. The
identity is realised at the canonical anchoring through the
specific Hubble closure
$h\!=\!\alpha_{\xi}s_{\rm face}N_{\rm gen}\!=\!27/40$ and
relies on the $(4,3)$ integer coincidences; it provides a
three-route consistency check on the Hubble closure that
cross-validates each of its three additive summands).
\item $Y_{p}\;-\;V_{us}\;=\;2\gamma^{2}\;=\;1/50$ \quad
(BBN primordial-helium fraction
$Y_{p}\!=\!(d{+}N_{\rm gen})^{2}\gamma^{2}/2\!=\!49/200$~\cite{Paper4B}
and Cabibbo angle
$V_{us}\!=\!\alpha_{\xi}s_{\rm face}\!=\!9/40$~\cite{Paper3}
differ by exactly twice the universal-leakage scale
$\gamma^{2}$. Three structurally independent sectors --- BBN
nucleosynthesis, $CKM$ quark mixing, and the universal
carrier-defect rate --- are linked through a single small
rational: $49/200\!-\!45/200\!=\!4/200\!=\!2\gamma^{2}$).
\item $H_{0}/(100\,{\rm km/s/Mpc})\;-\;\sin^{2}\theta_{23}^{\rm PMNS}\;=\;\gamma\;=\;1/10$ \quad
(dimensionless Hubble parameter and the atmospheric PMNS
mixing angle
$\sin^{2}\theta_{23}^{\rm NO}\!=\!1/2\!+\!\alpha_{\xi}/(d\,N_{\rm gen})\!=\!23/40$~\cite{Paper3}
differ by exactly the universal-leakage primitive
$\gamma$: $27/40\!-\!23/40\!=\!1/10$. Hubble cosmology
(P4B) and neutrino oscillations (P3) are interlocked at the
single-$\gamma$ offset level).
\item $n_{s}\;-\;\rho(k, k_{\Delta})\;=\;\varepsilon_{\rm sync}^{2}\;=\;1/200$ \quad
(inflation scalar spectral index
$n_{s}\!=\!193/200$~\cite{Paper4B} and carrier-axis-bundle
correlation
$\rho(k, k_{\Delta})\!=\!24/25\!=\!192/200$
(par.~\ref{par:cross_closure_identities}, owning paper P4)
differ by exactly the fluctuation--dissipation sync-defect
$\varepsilon_{\rm sync}^{2}\!=\!\gamma/2\!=\!1/200$. The
inflationary observable, the carrier graph statistic, and the
sync-channel primitive --- three structurally independent
quantities --- are linked through a single $\varepsilon^{2}$
offset).
\item $\Omega_{m}\;-\;\sin^{2}\theta_{W}\;=\;2d\,\gamma^{2}\;=\;8\gamma^{2}\;=\;2/25$ \quad
(matter-density fraction
$\Omega_{m}\!=\!47/150$~\cite{Paper4B} and electroweak mixing
angle
$\sin^{2}\theta_{W}\!=\!7/30\!=\!35/150$~\cite{Paper3} differ
by exactly $2d$ times the universal-leakage scale $\gamma^{2}$:
$47/150\!-\!35/150\!=\!12/150\!=\!2/25\!=\!8\gamma^{2}$. The
cosmological matter budget and the gauge-mixing angle are
interlocked at the order-$\gamma^{2}$ level by the
chirality-doubled-spatial-dimension prefactor $2d$).
\item $\tfrac{1}{4}\;-\;Y_{p}\;=\;\varepsilon_{\rm sync}^{2}\;=\;1/200$ \quad
(Bekenstein--Hawking entropy face $s_{\rm face}\!=\!1/d\!=\!1/4$
~\cite{Paper3}~L7 and BBN primordial-helium fraction
$Y_{p}\!=\!49/200$~\cite{Paper4B} differ by exactly the
sync-defect $\varepsilon_{\rm sync}^{2}$: $50/200\!-\!49/200\!=\!1/200$.
Black-hole entropy and primordial-helium synthesis are linked
at the sync-defect level, the same $\varepsilon^{2}$ that
appears in the $n_{s}\!-\!\rho$ offset above and that closes
the $H_{0}$ three-summand identity).
\item $\Omega_{\Lambda}\;+\;\sin^{2}\theta_{W}\;-\;Y_{p}\;=\;H_{0}/(100\,{\rm km/s/Mpc})\;=\;27/40$ \quad
(dark-energy fraction
$\Omega_{\Lambda}\!=\!103/150$~\cite{Paper4B}, electroweak
mixing angle $\sin^{2}\theta_{W}\!=\!7/30$~\cite{Paper3}, and
BBN primordial-helium fraction
$Y_{p}\!=\!49/200$~\cite{Paper4B} combine to the dimensionless
Hubble parameter:
$412/600\!+\!140/600\!-\!147/600\!=\!405/600\!=\!27/40$. Three
empirically and theoretically independent sectors --- late-time
dark-energy density, electroweak gauge mixing, and BBN
nucleosynthesis --- close on the cosmological Hubble parameter
through an exact small-integer rational combination at the
$(4, 3)$ anchoring point).
\item $Y_{p}\;+\;\sin^{2}\theta_{23}^{\rm PMNS}\;-\;\sigma_{8}^{2}\;=\;\gamma^{2}\;=\;1/100$ \quad
(BBN primordial-helium fraction $Y_{p}\!=\!49/200$,
atmospheric PMNS angle $\sin^{2}\theta_{23}\!=\!23/40$, and
matter-power amplitude squared $\sigma_{8}^{2}\!=\!\alpha_{\xi}^{2}\!=\!81/100$
sum to the universal-leakage scale squared:
$49/200\!+\!115/200\!-\!162/200\!=\!2/200\!=\!1/100\!=\!\gamma^{2}$.
The BBN, neutrino-oscillation, and structure-formation sectors
collapse to the carrier-defect rate squared).
\end{itemize}
These identities are \emph{not} universal in
$(d, N_{\rm gen})$: each relies on a specific integer
coincidence at the canonical anchoring (e.g.\
$N_{\rm gen}\!=\!d\!-\!1$, $2^{d}\!=\!d^{2}\!=\!16$,
$2N_{\rm gen}\!-\!1\!=\!d\!+\!1$). Their simultaneous
realisation is itself a structural feature of the
$(4,3)$ choice and a registered cross-corpus consistency
check: the cosmological, electroweak, $(\mathrm{SG})$-spectral
and curvature sectors all interconnect through small-integer
rationals at the empirically observed integer pair.
\fi
% END SANDBOX-MOVED cross-closure-identities block.

\section{Three load-bearing claims}
\label{sec:pillars}

The technical-foundations paper P6~\cite{Paper6} states the
load-bearing classification explicitly in its
``three claims are load-bearing'' paragraph: the
construction stands on
\begin{enumerate}
\item the \emph{bounded carrier with finite-$N$
Einstein-residual closure} of P4~\cite{Paper4} (emergent
gravity, lattice $\Lambda_{\rm lat}\!=\!19/15$ identity, four-point
Continuum Backreaction Identity);
\item the \emph{System-$\Rsys$ coefficient readout via the
bounded-operator construction} of P2~\cite{Paper2} (ten causal-wave
landings, formula-search audit, the four-projector spectral
uniqueness theorem);
\item the \emph{topology-labeled loop-class library mapping}
of P3~\cite{Paper3} (loop-class protocol, Yukawa-cluster
PROSP-01/04, Bekenstein--Hawking $S/A\!=\!1/4$).
\end{enumerate}
All other cross-sector closures ---
$\alpha_{\rm EM}$, $v_{\rm EW}$, $S_{\rm BH}/A\!=\!1/4$,
$n_{s}$, $A_{s}$, halo phenomenology, dark sector --- are
\emph{downstream} of the three load-bearing claims in the sense
defined in Section~\ref{sec:grammar}.

\section{Two branches: vacuum and matter-side chirality flip}
\label{sec:branches}

The bridge note~\cite{BridgeNote} formulates the corpus-wide
two-branch structure as a single \emph{loop-topological carrier}
operating in two distinct relational regimes:

\paragraph{Vacuum branch.} Pre-chirality-flip configuration
($\theta_{\rm chir}\!=\!0$, canonical
$(\axi,\gamma)\!=\!(9/10,1/10)$). Governs the electroweak-scale
closure ($v_{\rm EW}\!=\!246.2186$~GeV via P1~\cite{Paper1}'s
four-path consensus), the inflation-era and recombination-era
observables ($n_{s}, A_{s}, z_{*}, \theta_{*}$), and the
late-time dark-energy equation of state at leading order
($w_{\rm DE}\!=\!-39/40$).

\paragraph{Matter-side chirality-flip branch.} Post-$\pi/4$
matter-side branch with asymptotic/inverted anchor
$(\axi,\gamma)\!=\!(1/10,9/10)$: the chirality-flip boundary
itself sits at $\theta_{\rm chir}\!=\!\pi/4$ (where the two
coefficients cross), and the asymptotic post-flip anchor is
the inverted tuple $(1/10,9/10)$ approached as
$\theta_{\rm chir}$ continues past $\pi/4$. This branch governs
the matter-core phenomenology (top-$5\%$ K/Q closure), the
chirality-running $w_{\rm DE}(\theta_{\rm chir})$ at
next-to-leading order (P4B~\cite{Paper4B}), and the absolute
neutrino-mass identification $m_{1}\!=\!\gamma^{2}m_{3}$ via
the loop-class companion's Pure-Self-Energy resummation
(P3~\cite{Paper3}).

The two-branch decomposition is itself System-$\Rsys$-derived
rather than ad hoc: the chirality angle $\theta_{\rm chir}(N)$
runs algebraically with lattice resolution $N$~\cite{Paper4}
(the running formula is the one labelled
\emph{theta-running-first-principles} in P4), crossing
$\pi/4$ at a specific lattice resolution at which the
back-channel-projection numerical hierarchy inverts.

\section{Corpus map: which paper does what}
\label{sec:corpus-map}

\begin{table}[!htbp]
\centering
\small
\renewcommand{\arraystretch}{1.25}
\begin{tabular}{p{0.13\textwidth} p{0.36\textwidth} p{0.42\textwidth}}
\toprule
Paper & Title (short) & Role in the corpus \\
\midrule
P0 (this doc.) & Reader's guide & Entry-point, claim grammar,
                                  corpus map, falsifiers \\
P1~\cite{Paper1} & Hierarchy-bridge electroweak-scale closure
                                & $v_{\rm EW}\!=\!246.2186$~GeV
                                  four-path consensus,
                                  $|g|_{\overline{\rm MS}}\!=\!1.4187$,
                                  HBR + $\theta_{em}\!=\!\pi$
                                  load-bearing input to P4 \\
P2~\cite{Paper2} & Causal-wave landings + System-$\Rsys$ readout
                                & \textbf{Pillar (ii)}:
                                  bounded-operator construction;
                                  ten landings; four-projector
                                  spectral-uniqueness theorem;
                                  formula-search audit
                                  47/$\sim\!174$k \\
P3~\cite{Paper3} & Loop-class library + closure protocol
                                & \textbf{Pillar (iii)}: loop-class
                                  topology-labelled mapping;
                                  Bekenstein--Hawking
                                  $S/A\!=\!1/4$; SYE
                                  Yukawa/CKM/PMNS pipeline;
                                  PROSP-01/04 prospective registry \\
P4~\cite{Paper4} & Emergent-gravity bulk Einstein closure
                                & \textbf{Pillar (i)}: bounded
                                  carrier; $\Lambda_{\rm lat}^{\infty}\!=\!19/15$;
                                  Continuum Backreaction Identity;
                                  closure-domain ladder $N\!\in\![50,512]$ \\
P4A~\cite{Paper4A} & Schwarzschild + PPN companion
                                & Companion test:
                                  weak-field PPN parameters;
                                  $\bar\theta_{\rm QCD}$ vanishing \\
P4B~\cite{Paper4B} & Anisotropic source $\to$ DM/DE companion
                                & Companion: $T_{\mu\nu}^{\Xi}$
                                  decomposition; $w_{\rm DE}\!=\!-39/40$;
                                  chirality-running $w(z)$; 9-layer
                                  cosmological-constant dressing \\
P4C~\cite{Paper4C} & H3c witnesses companion
                                & Companion diagnostics for the
                                  Einstein-residual closure \\
P4D~\cite{Paper4D} & Atiyah--Singer chirality index companion
                                & Companion: integer-quantised
                                  vortex-DM winding number;
                                  topological-DM identification \\
Bridge note (appendix)~\cite{BridgeNote}
                  & P1$\leftrightarrow$P4 cross-file consistency
                                & Lemma-level: common
                                  loop-topological carrier;
                                  version-lock matrix;
                                  cross-paper $\Rsys$-tuple
                                  consistency \\
P6~\cite{Paper6} & Relational-UV-closure technical foundations
                                & Master technical document:
                                  $H_{UV}$, $S_{UV}$, $Z_{UV}$,
                                  branch construction, closure
                                  ledger, load-bearing-vs-downstream
                                  registry \\
\bottomrule
\end{tabular}
\caption{Corpus map. P0 (this document) is the reader's guide
and entry point; P2, P3, P4 are the three load-bearing claims
of Section~\ref{sec:pillars}; P1, P4A--D, and P6 supply
supporting structure or technical foundations; the bridge
note is a cross-file consistency appendix.}
\label{tab:corpus-map}
\end{table}

\section{Claim grammar}
\label{sec:grammar}

The corpus uses a registered vocabulary of closure tiers and
claim classes. Each is defined formally in one of the technical
papers; the present section is a routing table only.

\paragraph{Closure tiers (precision against external anchor).}
\begin{itemize}
\item \textbf{EXACT}: relative residual at or below the
paper-specific EXACT cut, corresponding to a System-$\Rsys$
identity that closes at machine precision under the rational
reduction $\Rsys$ (see P6 §1 definition).
\item \textbf{PRECISE}: relative residual $\le 2.5\%$; closes
the framework prediction within standard-fit experimental
uncertainty.
\item \textbf{FACTOR2}: closure form correct but numerical value
within a factor of 2; reported as a structural-form
identification, not a precision claim.
\item \textbf{STRUCTURAL\_CLOSED}: closure form derived
algebraically; numerical match secondary
(e.g.\ $\Sigma_{f}Y_{f}^{2}\!=\!10$).
\end{itemize}

\paragraph{Terminological disclaimer on the word ``EXACT''.}
Throughout this corpus the labels EXACT, PRECISE and FACTOR2
are \emph{band labels}, not claims of mathematical equality.
EXACT here means ``the framework value lies inside the
indicated relative-residual band against the external
anchor''; a 1\% residual against an experimental measurement
is not exact in the mathematical sense of $a = b$ in
$\mathbb Q$. Where the corpus does report a true algebraic
identity in $\mathbb Q$ (e.g.\ $\alpha_\xi/2 - 2\gamma = 1/4$
for $S_{\rm BH}/A$, or
$\Sigma_{f}Y_{f}^{2}\!=\!10$ for the
hypercharge-squared sum rule), the closure is flagged as
\textbf{STRUCTURAL\_CLOSED} or written with the explicit
identity symbol $\equiv$ in the body text. Readers comparing
the corpus against algebraic-identity theories (e.g.\
representation-theoretic constructions in $\mathbb Q$) should
treat EXACT-band labels as ``Tier-1 closure within band'',
not as algebraic identities.

We use Tier-1 / Tier-2 / Tier-3 interchangeably with EXACT /
PRECISE / FACTOR2; the Tier-$k$ vocabulary is preferred when
the rhetoric of EXACT could mislead an external reader who
expects mathematical equality.

\paragraph{Paper-specific EXACT cuts.} The corpus uses
three EXACT cuts depending on the observable class. The
present routing table makes them explicit so a reader
cross-comparing two papers does not confuse a stricter
EXACT-label in one paper with a looser one in another. A
claim that is EXACT under a stricter cut is automatically
EXACT under any looser cut (the implication runs from
left to right).
\begin{center}
\begin{tabular}{@{}p{0.27\textwidth} p{0.13\textwidth} p{0.48\textwidth}@{}}
\toprule
Paper / observable class & EXACT cut & Rationale \\
\midrule
P2 (causal-wave landings) & $|r| \le 0.01\%$ & EXACT-target
algebraic landings; sub-experimental-uncertainty closure \\
P3 (loop-class library) & $|r| < 0.4\%$ & loop-corrected
observables compared against PDG/NuFIT measurement bands \\
P0 / P4 / P4A--D / P6 / bridge & $|r| \le 1\%$ & global
manifest band covering all observable classes; superset of
P2 and P3 cuts \\
\bottomrule
\end{tabular}
\end{center}
The PRECISE cut is uniform across the corpus
($|r| \le 2.5\%$) and the FACTOR2 cut is uniform
($|r| \le 100\%$, i.e.\ within a factor of two).
A claim that closes at $|r| \le 0.01\%$ is EXACT under any
of the three cuts; a claim at $|r| \in (0.01\%, 0.4\%]$ is
EXACT under P3, P0 and the global manifest but downgrades to
PRECISE under P2; a claim at $|r| \in (0.4\%, 1\%]$ is EXACT
under P0 and the global manifest but PRECISE under P3.

\paragraph{Claim classes (load-bearing classification).}
The unified manifest (file
\path|data/corpus_claims_manifest.json| in P6's
repository~\cite{Paper6}) registers every main claim under
one of five labels:
\begin{itemize}
\item \textbf{Load-bearing.} The construction stands on this
claim. Three exist: pillars (i)--(iii) of
Section~\ref{sec:pillars}.
\item \textbf{Downstream.} An algebraic or audit consequence
of the same System-$\Rsys$ tuple read off by the load-bearing
claims (e.g.\ $\alpha_{\rm EM}\!=\!\gamma^{2}\axi^{N_{\rm gen}}$,
$\sigma_{8}\!=\!\axi^{2}$,
$n_{s}\!=\!1\!-\!7\gamma^{2}/2$). The vast majority of the
closures of Table~\ref{tab:top-closures} fall here.
\item \textbf{Supporting witness.} An independent audit
diagnostic that does not extend the closure surface but
confirms it (e.g.\ P4C's H3c gravitational-coupling-scaling
diagnostics).
\item \textbf{External contact.} A point of contact with the
existing literature (Pauli--Zeldovich cancellation,
Schwinger--Dyson resummation, Reed--Simon VIII.5).
\item \textbf{Open target.} A registered observable whose
closure is not yet established within the corpus and which is
flagged as the framework's primary near-term theoretical
target.
\end{itemize}

\section{Main closures at a glance}
\label{sec:top-closures}

The list below is a curated subset of the $\sim\!50$ registered
closures of the corpus, selected for highest external-anchor
visibility. Each entry is owned by the cited paper; this table
re-states them in one place.

\setlength{\tabcolsep}{3pt}
\renewcommand{\arraystretch}{1.18}
\begin{longtable}{@{}p{0.17\textwidth} p{0.25\textwidth} p{0.18\textwidth} p{0.16\textwidth} p{0.07\textwidth}@{}}
\caption{Top-tier corpus closures, one row per observable.
``Tier'' is the precision class against the external anchor:
\textbf{EXACT} ($\equiv$\,Tier-1; $|r| \le 1\%$),
\textbf{PRECISE} ($\equiv$\,Tier-2; $|r| \le 2.5\%$),
\textbf{FACTOR2} ($\equiv$\,Tier-3; within $2\times$). All three
are band labels, not claims of mathematical equality; see the
terminological disclaimer in Sec.~\ref{sec:grammar}.
``Class'' is the load-bearing classification of
Section~\ref{sec:grammar}. Rows marked \emph{load-bearing} are
components of the three load-bearing blocks
(pillars~(i)--(iii) of Section~\ref{sec:pillars}), not
additional independent pillars; the corpus has \emph{exactly
three} load-bearing pillars and the multiplicity of
load-bearing rows in the table reflects the multiple
observables that each pillar supplies.}
\label{tab:top-closures}\\
\toprule
Observable & System-$\Rsys$ form & Tier & Class & Owner \\
\midrule
\endfirsthead

\toprule
Observable & System-$\Rsys$ form & Tier & Class & Owner \\
\midrule
\endhead

\multicolumn{5}{c}{\textit{Cosmology}} \\
$H_{0}$ & $100\,\axi\,s_{\rm face}\,N_{\rm gen}\!=\!67.5$ km/s/Mpc & EXACT & downstream & P4B \\
$\sigma_{8}$ & $\axi^{2}\!=\!81/100$ & PRECISE & downstream & P4B \\
$S_{8}$ & $\axi^{2}\sqrt{\Omega_{m}/0.3}\!=\!0.828$ & EXACT & downstream & P4B \\
$\Omega_{m}$ & $47/150$ & PRECISE & downstream & P4B \\
$\Omega_{\Lambda}$ & $103/150$ & PRECISE & downstream & P4B \\
$\Omega_{\rm DM}h^{2}$ & $243/2000\!=\!0.1215$ & PRECISE & downstream & P4B \\
$w_{\rm DE}$ & $-1+\varepsilon_{\rm sync}^{4}/\gamma\!=\!-39/40$ & PRECISE & downstream & P4B \\
$n_{s}$ & $1-7\gamma^{2}/2$ & EXACT & downstream & P4B \\
$A_{s}$ & $21\gamma^{10}$ & EXACT & downstream & P4B \\
$z_{*}$ & $(2d{+}1)(2d{+}N_{\rm gen})^{2}\!+\!{\rm sub}$ & EXACT & downstream & P4B \\
$Y_{p}$ & $49\gamma^{2}/2$ & EXACT & downstream & P4B \\
$\Sigma m_{\nu}$ (NH-min, dressed) & $m_{2}^{\rm dr}\!+\!m_{3}^{\rm dr}\!=\!0.0588$ eV; seesaw refinement $\gamma^{2}m_{3}\!+\!m_{2}\!+\!m_{3}\!\simeq\!0.0593$ eV & EXACT & downstream & P4B \\
$m_{3}^{(\nu)}$ atmospheric & $m_{e}\gamma^{d+3}(1\!-\!\gamma^{2}(d{+}N_{\rm gen})/4)\!=\!0.0502$ eV & EXACT & downstream & P4B \\
$m_{2}^{(\nu)}$ solar & $m_{e}\gamma^{d+4}(\pi/2)(1\!+\!\gamma^{2}(d{+}N_{\rm gen}))\!=\!0.00859$ eV & EXACT & downstream & P4B \\
$\eta_{B}$ & $(d{+}N_{\rm gen})^{2}\gamma^{10}/(2d)\!=\!49/(8\!\cdot\!10^{10})$ & PRECISE & downstream & P4B \\
$D/H$ & $(d{+}1)\gamma^{5}/2\,(1\!+\!\gamma^{2}(d{+}N_{\rm gen})/4)\!=\!2.544\!\times\!10^{-5}$ & EXACT & downstream & P3 \\
\midrule
\multicolumn{5}{c}{\textit{Electroweak \& Standard Model}} \\
$v_{\rm EW}$ & 4-path consensus $246.2186$ GeV & PRECISE & load-bearing input & P1 \\
$\alpha_{\rm EM}$ & $\gamma^{2}\axi^{N_{\rm gen}}\!=\!729/100{,}000$ & EXACT & downstream & P3 \\
$\sin^{2}\theta_{W}$ & $1/d - \gamma/(2N_{\rm gen})\!=\!7/30$ & PRECISE & downstream & P3 \\
$V_{us}$ & $\axi\,s_{\rm face}\!=\!9/40$ & EXACT & downstream & P3 \\
$m_{t}$ & spectral Yukawa eigenvalue $\to 174.1$ GeV & EXACT & downstream & P3 \\
$m_{\mu}/m_{e}$ & $400 N_{\rm gen}(d{+}1)/P_{29}\!=\!6000/29\!=\!206.897$ & EXACT & downstream & P3 \\
$a_{\mu}^{\rm Schwinger}$ & $\alpha_{\rm EM}/(2\pi)$ & PRECISE & downstream & P3 \\
\midrule
\multicolumn{5}{c}{\textit{Gravity \& black-hole sector}} \\
$\Lambda_{\rm lat}^{\infty}$ & $3\axi/2+\esync-\gamma(N_{\rm gen}{+}1)/N_{\rm gen}\!=\!19/15$ & EXACT & load-bearing & P4 \\
$\Lambda_{t}^{\infty}$ & $\axi^{2}\!=\!81/100$ & EXACT & load-bearing & P4 \\
$\Lambda_{s}^{\infty}$ & $-\gamma^{2}/2\!=\!-1/200$ & EXACT & load-bearing & P4 \\
$S_{\rm BH}/A$ & $s_{\rm face}\!=\!1/d\!=\!1/4$ & EXACT & load-bearing & P3 \\
PPN $\gamma_{\rm PPN}$ & induced from Einstein closure & PRECISE & downstream & P4A \\
$\rho_{\Lambda}^{\rm pred}/\rho_{\Lambda}^{\rm obs}$ & 9-layer dressing $\to 1.001$ & EXACT & downstream & P4B \\
\midrule
\multicolumn{5}{c}{\textit{Dark sector \& halo phenomenology}} \\
$g_{\dagger}$ (RAR) & $cH_{0}/(2\pi\axi)$ & PRECISE & downstream & P4B \\
BTFR $A$ & $\sim\!(G_{N}H_{0})^{-1}$ & PRECISE & downstream & P4B \\
NFW vs Burkert branching & polarity-violation rate & STRUCTURAL & downstream & P4B \\
SPARC pass-fraction & $79/129$ within domain $\mathcal{D}_{\rm halo}$ & PRECISE & downstream & P4B \\
Vortex-DM $w_{\rm DM}$ & $0$ (integer-quantised winding) & STRUCTURAL & downstream & P4D \\
\midrule
\multicolumn{5}{c}{\textit{Foundations}} \\
4-projector basis & Reed--Simon VIII.5 spectral-uniqueness & EXACT & load-bearing & P2 \\
10 causal-wave landings & 47 forms / $\sim\!174$k searched & EXACT & load-bearing & P2 \\
Loop-class library & topology-labelled protocol & STRUCTURAL & load-bearing & P3 \\
$\sum_{f}Y_{f}^{2}\!=\!10$ & three-generation chiral content & EXACT & downstream & P6 \\
$\bar\theta_{\rm QCD}\!=\!0$ & charge-neutral vortex/antivortex condensate (P4 + P4D + P3 chirality-pair) & STRUCTURAL & downstream & P4D \\
Emergent causal time / $3{+}1$ Lorentz layer & causal partial order on $\Xi$-paths $+$ massless transport light cone; $d_{\rm spatial}\!=\!3$, $d_{\rm spacetime}\!=\!4$, $f_{\rm LC}\!=\!\log_{2}(9/8)\!=\!0.169925$ & STRUCTURAL & supporting & P4 \\
Emergent UV time $t$ (Page--Wootters) & $t$ = Fiedler-mode-phase clock eigenvalue of $\Xi$-Laplacian; $\hat C|\Psi\rangle\!=\!0$ Wheeler--DeWitt constraint & STRUCTURAL & supporting & P6 \\
\bottomrule
\end{longtable}
\setlength{\tabcolsep}{6pt}
\renewcommand{\arraystretch}{1.0}

\section{Falsification handles}
\label{sec:falsifiers}

The framework registers explicit falsification triggers
against current and near-term cosmological and laboratory
data. Closure of these handles within the next $\sim\!2$ years
is the principal external test of the framework.

\paragraph{Cross-reference to owning-paper falsifier IDs.}
The six corpus-wide falsifiers below map to the
owning-paper local falsifier labels as registered in each
companion paper's \texttt{Falsification triggers} section.
\begin{center}
\small
\begin{tabular}{@{}p{0.40\textwidth} p{0.28\textwidth} p{0.22\textwidth}@{}}
\toprule
P0 falsifier (this section) & Owning paper & Local label \\
\midrule
1. $\Sigma m_{\nu}$ vs DESI / CMB-S4 / Roman / Euclid
   & P3, P6        & P3 FAL-$\Sigma m_{\nu}$, P6 §FAL-cosmo \\
2. Dynamical $w(z)$ vs DESI BAO running
   & P4B           & P4B FAL-$w_{a}$ \\
3. Hubble tension on non-Cepheid distance ladders
   & P4B           & P4B FAL-$H_{0}$ \\
4. Tensor-to-scalar ratio $r$ vs CMB-S4
   & P4B           & P4B FAL-$r$ \\
5. Particle-DM direct detection (LZ / XENON-nT / PandaX / DarkSide)
   & P4D           & P4D F1, F2, F3 \\
6. Halo regime signatures (microlensing, EDGES 21-cm)
   & P4, P4B       & P4B FAL-halo \\
\bottomrule
\end{tabular}
\end{center}
A reader following one of these falsifiers across to its
owning paper finds the same numerical thresholds at higher
technical detail and the reproducer pointers.

\paragraph{Cosmological-data-decidable within $\sim\!2$ years.}
\begin{enumerate}
\item \emph{$\Sigma m_{\nu}$ vs DESI / CMB-S4 / Roman / Euclid.}
The framework predicts $\Sigma m_{\nu}\!\simeq\!0.060$~eV via
$m_{1}\!=\!\gamma^{2}m_{3}$~\cite{Paper3,Paper6}. The current
DESI DR2 BAO + DR1 full-shape bound is
$\Sigma m_{\nu}\!<\!0.0642$~eV~\cite{DESI2025}; the framework
prediction sits $6.1\%$ below this. A confirmed tightening
below the NH oscillation minimum $\simeq\!0.0586$~eV would
falsify the $\gamma^{2}m_{3}$ identification.
\item \emph{Dynamical $w(z)$ vs DESI BAO running.} The framework
predicts $|w_{a}|\!\lesssim\!\gamma^{2}\!=\!0.01$
from the chirality-running structure
$w_{\rm DE}(\theta_{\rm chir})$~\cite{Paper4B}. The DESI 2024+2025
central $w_{a}\!\approx\!-0.4$ at $\sim\!3\sigma$ lies
$\sim\!40\!\times$ outside the framework band; consolidation
to $>\!5\sigma$ in DESI-DR3 would falsify the
$\varepsilon_{\rm sync}^{4}/\gamma$ closure.
\item \emph{Hubble tension on non-Cepheid distance ladders.}
The framework predicts $H_{0}\!=\!67.5$~km/s/Mpc independent
of distance-ladder method~\cite{Paper4B}. Confirmed
$H_{0}\!\in\![72,74]$~km/s/Mpc on a non-Cepheid local ladder
(gravitational-wave sirens, TDCOSMO, JWST-anchored TRGB at
$\sigma_{H_{0}}\!\lesssim\!1$ km/s/Mpc) would falsify the
distance-ladder-systematic attribution of the SH0ES
$73.04\!\pm\!1.04$ discrepancy.
\item \emph{Tensor-to-scalar ratio $r$ vs CMB-S4.} The
framework predicts $r\!\sim\!\gamma^{3}(d{+}N_{\rm gen})\!\simeq\!7\!\times\!10^{-3}$
within the end-of-2025 Planck--SPT--ACT--BICEP/Keck--DESI 95\% bound
$r\!<\!0.034$~\cite{Balkenhol2025} (subsuming the earlier BK18 bound
$r_{0.05}\!<\!0.036$~\cite{BICEPKeck2021}). CMB-S4 sensitivity
$\sigma(r)\!\sim\!10^{-3}$ within $\sim\!3$ years will provide
a definitive detection or upper-bound test.
\end{enumerate}

\paragraph{Laboratory-data-decidable.}
\begin{enumerate}
\setcounter{enumi}{4}
\item \emph{Particle-DM direct detection.} The framework
identifies dark matter with a relational vortex-defect network
of integer-quantised $U(1)$
winding~\cite{Paper4D}; this is topological, not particulate. The
framework therefore predicts \emph{null} results on LZ
($\sigma_{\rm SI}\!<\!2.2\!\times\!10^{-48}$\,cm$^{2}$ at
$40$ GeV/c$^{2}$~\cite{LZ2024}), XENON-nT~\cite{XENONnT2023},
PandaX-4T~\cite{PandaX2024}, and DarkSide~\cite{DarkSide2023}
down to the neutrino floor. A positive nuclear-recoil signal
at any future facility would falsify the topological-DM
identification.
\item \emph{Halo regime signatures.} Microlensing of
vortex-network clumps at the chirality-flip mass scale
$M_{\rm flip}$~\cite{Paper4} should show a discrete
event-frequency distribution rather than a continuous PBH
spectrum; 21-cm absorption-depth deviations from $\Lambda$CDM
in the EDGES band correlate with the vortex-density power
spectrum at $z\!\sim\!17$.
\end{enumerate}

\paragraph{Pre-registered prospective targets.}
The loop-class companion paper P3~\cite{Paper3} registers a
prospective cluster registry --- the file
\path|prospective_cluster_registry.json|
in the loop-class repo --- that pre-registers three
loop-class observables for future $\Rsys$-falsification
opportunities at the P3 protocol level.

\paragraph{Composition theorem conditionality: the
$(\mathrm{SG})$ axiom.}
The composition theorem of P6~\cite{Paper6} states that
$(M1)$, $(M2)$, $(M4)$ are unconditional theorems of the cited
papers, while $(M3)$ --- the emergent-Einstein closure --- is
\emph{conditionally rigorous} on three inputs: (i) the eight
admissibility conditions $A_{1}$--$A_{8}$ on the carrier
sequence, (ii) the discharge-table implication
$(A_{5})\!+\!(A_{6})\!\Rightarrow\!CD(K_{\mathrm{CD}},N)$ open at
theorem level in P4, and (iii) the empirical Symanzik asymptotes
underlying the Continuum Backreaction Identity
($T_{00}^{\Xi}$ leg $R^{2}\!=\!0.54$, auxiliary
$G_{00,\mathrm{med}}\!\to\!0$ at $R^{2}\!=\!0.99$). The P6
sharpening reduces (i) to a \emph{single} structural axiom:
the uniform spectral gap $(\mathrm{SG})$ on the $\xi_{N}$-weighted
normalised graph Laplacian, with conjectured asymptote
$\lambda_{\infty}\!=\!3/8\!=\!(d{-}1)/(2d)$. The axiom is
\emph{empirically certified} ($10$-regime ladder, $184$ seeds,
Symanzik-$1$ asymptote $0.3789$ at $1.0\%$ relative residual to
$3/8$) but \emph{not analytically discharged}. A GPU
decomposition study (P4 Prop.\ Statistical-fingerprint, 2026)
further reduces the open analytical question to a concrete
graph-theoretic statement: the spectral gap of the
Newman-clustered ensemble with the carrier's empirical
(degree, triangle) distribution equals $7/24$ at $+0.36\%$ relative
residual --- and \emph{only} this ensemble reproduces the
carrier asymptote (Erd\H{o}s--R\'enyi, configuration model,
random-regular, and Watts--Strogatz all miss by $-40$ to $+66\%$).
The corresponding analytical tool is the cavity / resolvent
method for sparse networks with triangles; the remaining
Step-$4$ deliverable is therefore concretely the closed-form
cavity fixed-point solution for the carrier's clustered-degree
distribution. \emph{Falsification handle}: any rigorous proof
that the carrier-clustered-graph cavity solution does \emph{not}
yield $7/24$, or any tightened empirical measurement that excludes
$3/8$ on a $\Rsys$-aware ladder at $\geq\!5\sigma$, would close
the axiom against the framework rather than for it.

\paragraph{Physical mechanism of the corpus closures.}
\label{par:mechanism_status_cross_projection}
Each closure of the corpus is classified along two independent
axes: a \emph{numerical-precision tier} (EXACT / PRECISE /
FACTOR2; the closure-grammar registry above) and a
\emph{physical-mechanism status} (rigorous derivation, physical
plausibility argument, or open). Every numerical landing in
the corpus comes with either a rigorous physical derivation or
a concrete plausibility argument connecting the System-$\Rsys$
form to the observable; no closure is registered as
purely numerical without physical reading.

A subset of the closures admit no single-route derivation but
are nonetheless established by a \emph{cross-projection
argument}: two or more independent in-corpus chains connect to
the observable through standard physics, and their numerical
predictions converge at the PRECISE tier. The three cosmological
and geometric closures listed below are of this kind.

\begin{itemize}
\item \emph{Local Hubble constant
$H_{0}\!=\!67.5$~km/s/Mpc}. The dimensionless reading
$h\!:=\!H_{0}/(100\,\mathrm{km/s/Mpc})\!=\!
\alpha_{\xi}\,s_{\rm face}\,N_{\rm gen}\!=\!27/40$ from the
carrier-side product is also delivered by the standard
Friedmann equation at $z\!=\!0$, with the framework's
independent dressing of $\rho_{\Lambda}$~\cite{Paper4B} and
its closure $\Omega_{\Lambda}\!=\!1-\Omega_{m}\!=\!103/150$
as inputs. The carrier-side route gives $67.50$ and the
Friedmann route $66.93$; the two converge at $0.85\%$
relative residual, each within Planck~2018 $1\sigma$ of the
empirical $67.4\!\pm\!0.5$.
\item \emph{Lattice cosmological-constant asymptote
$\Lambda_{\rm lat}^{\infty}\!=\!19/15$}. The operator-side
decomposition $17/20\!+\!5/12$ (non-scalar Clifford-channel
reaction rate plus spinor-trace--generation correction,
\cite{Paper4}) is algebraically equivalent to the carrier-side
decomposition
$1\!+\!d/((d{+}1)N_{\rm gen})\!=\!1\!+\!4/15$
(trivial unit carrier-baseline plus per-generation
chirality-spin defect distributed over the sync-extended cone
bundle). Two independent physical projections of the same
limit, agreeing identically at $19/15$.
\item \emph{Curvature-dimension lower bound
$\mathrm{CD}(K_{\rm CD}\!\geq\!0, N)$}. Three independent
Bakry--\'Emery percentile statistics each match a
$\lambda_{\infty}\!\times\!(d, N_{\rm gen})$-rational ---
$K_{p_{1}}\!=\!\lambda_{\infty}(d{+}1)/d\!=\!15/32$ at
$0.19\%$ residual,
$K_{p_{5}}\!=\!\lambda_{\infty}N_{\rm gen}(d{+}N_{\rm gen})/2^{d}
\!=\!63/128$ at $0.04\%$ residual,
$K^{\rm mean}\!=\!\lambda_{\infty}(4d{+}1)/(2d{+}N_{\rm gen})
\!=\!51/88$ at $0.10\%$ residual. The three-fold convergence
identifies the Bakry--\'Emery curvature distribution as a
structural function of the $(\mathrm{SG})$-axiom
$\lambda_{\infty}$ and the integer inputs $(d, N_{\rm gen})$,
extending the qualitative Ollivier--Lin--Lu--Yau corollary
$K_{\rm CD}\!\geq\!0$ to a quantitative cross-projection.
\end{itemize}

The cross-projection arguments are not stand-alone derivations:
each identity rests on standard physics
(Friedmann; Bakry--\'Emery $\Gamma_{2}$-calculus;
algebraic-decomposition equivalence) plus the framework's
prior closures of the routes' inputs. The load-bearing
evidence is the empirical convergence of two or more
independent in-corpus chains at residuals of $0.04$--$0.85\%$,
all inside the bootstrap $95\%$ confidence intervals of the
corresponding empirical measurements. The single remaining
mechanism conditionality of the corpus is therefore the
discharge of the $(\mathrm{SG})$ axiom
$\lambda_{\infty}\!=\!3/8$ itself: a closed-form derivation
would, via the cross-projection chains above, simultaneously
deliver rigorous physical mechanisms for all three closures
of this kind.

\paragraph{Statistical multiplicity / look-elsewhere.}
The corpus reports many closures, which invites a
look-elsewhere objection. The defence does \emph{not} rest on
post-hoc $p$-value-inflation corrections; it rests on five
structural features that a random-search artefact would not
exhibit. \emph{(i)~Dependency graph}: the closures are not
independent fits but a directed graph rooted in the two
integers $(d, N_{\rm gen})$ --- every System-$\Rsys$ rational
is a fixed algebraic function of those two inputs, so the
effective number of free choices is two, not the number of
closures. \emph{(ii)~Frozen readouts}: the five carrier
coefficients are read off a bounded-operator construction
(P2~\cite{Paper2}) once and then frozen; no downstream closure
re-tunes them. \emph{(iii)~Negative controls}: the loop-class
protocol (P3~\cite{Paper3}) has explicit forbidden classes and
a closure domain outside which the protocol is \emph{required}
to fail, and several candidate identifications have been
falsified on those controls rather than retained.
\emph{(iv)~Prospective registries}: P3 pre-registers
loop-class observables (\path|prospective_cluster_registry.json|)
before their external values enter the closure, converting them
from fits into predictions. \emph{(v)~No target-fitting}: the
System-$\Rsys$ forms carry no fitted parameter --- the
admissible alphabet and the formula-search audit of
P2~\cite{Paper2} are disclosed, and a closure either lands on a
System-$\Rsys$ rational or it does not. A look-elsewhere
artefact would require all five to hold by coincidence; the
framework is therefore best contested through the registered
falsification handles above, not through a multiplicity
correction on a closure count.

\section{Where to read next}
\label{sec:routing}

The natural reading order depends on the reader's question:

\paragraph{``What is the theory?''.} Read this guide (P0)
first. Then read P6~\cite{Paper6} for the technical
foundations, P2~\cite{Paper2} for the System-$\Rsys$ readout,
P3~\cite{Paper3} for the loop-class protocol, and
P4~\cite{Paper4} for the main emergent-GR claim. The bridge
note~\cite{BridgeNote} is best read afterward as a
cross-file consistency appendix.

\paragraph{``Where do the System-$\Rsys$ coefficients come
from?''.} P2~\cite{Paper2} (causal-wave landings + formula-search
audit + spectral-uniqueness theorem).

\paragraph{``How does general relativity emerge?''.}
P4~\cite{Paper4} (bounded carrier + Einstein-residual closure +
Continuum Backreaction Identity); then P4A~\cite{Paper4A} for
the weak-field PPN companion.

\paragraph{``How do the Standard-Model masses and mixings
close?''.} P3~\cite{Paper3} (loop-class library + SYE
Yukawa-pipeline + $\alpha_{\rm EM}$-mediated identities).

\paragraph{``How do dark matter and dark energy fit in?''.}
P4B~\cite{Paper4B} (anisotropic source decomposition,
$w_{\rm DE}$, cosmological-constant dressing) and
P4D~\cite{Paper4D} (vortex-DM topological identification).

\paragraph{``What is the electroweak-scale closure?''.}
P1~\cite{Paper1} (HBR + $\theta_{em}\!=\!\pi$ four-path
consensus for $v_{\rm EW}$).

\paragraph{``Is the corpus internally consistent?''.} The
bridge note~\cite{BridgeNote} (technical-consistency proof
P1$\leftrightarrow$P4) and the unified manifest file
\path|relational-uv-closure-repro/data/corpus_claims_manifest.json|
of P6~\cite{Paper6}.

\section*{Closing note}

The relational carrier theory is presented as a candidate
closure construction whose external validation is the principal
next step. The corpus does not claim to be a final accepted
theory of nature; it claims that the load-bearing pillars
(i)--(iii) of Section~\ref{sec:pillars} are internally
consistent, that the downstream closures of
Table~\ref{tab:top-closures} follow algebraically from the
load-bearing pillars within the System-$\Rsys$ rational
reduction, and that the falsification handles of
Section~\ref{sec:falsifiers} provide a decidable near-term
test surface. Each individual claim is owned by the cited
paper; the present guide is an entry point and a map, nothing
more.

\section*{Data and software availability}
\label{sec:data-availability}
The reproducibility package
(\url{https://github.com/TerraSignum/relational-carrier-manifest}) bundles all
input data files (under \path|data/|), reproducer scripts
(\path|src/verify_*.py| and \pyscript{causal-wave-landings-repro}{src/make_figures.py}{make_figures}), and
unit tests (\path|tests/|). Cryptographic provenance is provided by
\path|data/SHA256SUMS| (GNU \texttt{sha256sum -c} compatible),
and the publication-prep frozen state is documented in
\path|RELEASE_v1_0_0.md| with full release notes in
\path|CHANGELOG.md|. Software metadata is in \path|pyproject.toml|
and \path|CITATION.cff| (Citation File Format 1.2). The package is
licensed under MIT; the manuscript text is licensed under
CC-BY-4.0.

\bibliographystyle{unsrt}
\begin{thebibliography}{99}

\bibitem{Paper1}
Bucciarelli, S.,
\emph{A defect-field two-loop correction to the
Hierarchy-Bridge-Relation electroweak scale}
(\path|hbr-ew-scale-repro|, 2026).

\bibitem{Paper2}
Bucciarelli, S.,
\emph{A five-coefficient causal-wave ansatz and ten cross-sector
benchmark closures}
(\path|causal-wave-landings-repro|, 2026).

\bibitem{Paper3}
Bucciarelli, S.,
\emph{Topology-labeled loop-class mappings for cross-sector
physical observables}
(\path|loop-class-closure-repro|, 2026).

\bibitem{Paper4}
Bucciarelli, S.,
\emph{Emergent Einstein dynamics from relational metric closure
on finite correlation lattices}
(\path|emergent-gr-closure-repro|, 2026).

\bibitem{Paper4A}
Bucciarelli, S.,
\emph{Schwarzschild defect, post-Newtonian limits, and the
black-hole sector from a relational emergent-metric
construction}
(\path|emergent-gr-schwarzschild-ppn-repro|, 2026).

\bibitem{Paper4B}
Bucciarelli, S.,
\emph{Anisotropic source tensor and the dark-matter /
dark-energy mixture decomposition on a relational lattice}
(\path|emergent-gr-anisotropic-source-dm-de-repro|, 2026).

\bibitem{Paper4C}
Bucciarelli, S.,
\emph{Multi-observable indirect-witness chain bounding the
Einstein-identity gap on a relational emergent-gravity lattice}
(\path|emergent-gr-h3c-witnesses-repro|, 2026).

\bibitem{Paper4D}
Bucciarelli, S.,
\emph{A discrete Atiyah--Singer-type bridge from the
chirality-balance witness to the Einstein-identity gap on the
relational $\Xi$-graph}
(\path|emergent-gr-atiyah-singer-chirality-repro|, 2026).

\bibitem{Paper6}
Bucciarelli, S.,
\emph{Full UV-closure construction of the relational
QFT--Einstein carrier: a two-phase loop-topological framework
with vacuum and matter branches}
(\path|relational-uv-closure-repro|, 2026).

\bibitem{BridgeNote}
Bucciarelli, S.,
\emph{A common loop-topological carrier for electroweak and
emergent-gravity sectors (Technical consistency note across
Papers~1--4)}
(\path|qft-einstein-bridge-repro|, 2026).

\bibitem{DESI2025}
W.~Elbers \textit{et al.} (DESI Collaboration),
\emph{Constraints on neutrino physics from DESI DR2 BAO and DR1
full shape},
Physical Review D \textbf{112}, 083513 (2025); arXiv:2503.14744.

\bibitem{BICEPKeck2021}
P.~A.~R.~Ade \textit{et al.} (BICEP/Keck Collaboration),
\emph{Improved constraints on primordial gravitational waves
using Planck, WMAP, and BICEP/Keck observations through the
2018 observing season} (BK18),
Physical Review Letters \textbf{127}, 151301 (2021); arXiv:2110.00483.

\bibitem{Balkenhol2025}
L.~Balkenhol \textit{et al.},
\emph{Inflation at the End of 2025: Constraints on $r$ and $n_s$
Using the Latest CMB and BAO Data},
arXiv:2512.10613 (December 2025).
Combined Planck--SPT--ACT--BICEP/Keck + DESI BAO analysis;
$n_{s}\!=\!0.9682\!\pm\!0.0032$ (CMB only),
$0.9728\!\pm\!0.0029$ with BAO; 95\% C.L. upper bound $r\!<\!0.034$.

\bibitem{LZ2024}
J.~Aalbers \textit{et al.} (LUX-ZEPLIN Collaboration),
\emph{Dark matter search results from 4.2 tonne-years of exposure
of the LUX-ZEPLIN (LZ) experiment},
Physical Review Letters \textbf{135}, 011802 (2025); arXiv:2410.17036.

\bibitem{XENONnT2023}
E.~Aprile \textit{et al.} (XENON Collaboration),
\emph{First dark matter search with nuclear recoils from the
XENONnT experiment},
Physical Review Letters \textbf{131}, 041003 (2023); arXiv:2303.14729.

\bibitem{PandaX2024}
Z.~Bo \textit{et al.} (PandaX Collaboration),
\emph{Dark matter search results from 1.54 tonne$\cdot$year exposure
of PandaX-4T},
Physical Review Letters \textbf{134}, 011805 (2025); arXiv:2408.00664.

\bibitem{DarkSide2023}
P.~Agnes \textit{et al.} (DarkSide-50 Collaboration),
\emph{Search for low-mass dark matter WIMPs with 12 ton-day exposure
of DarkSide-50},
Physical Review D \textbf{107}, 063001 (2023); arXiv:2207.11966.

\end{thebibliography}

\end{document}
